% SPDX-FileCopyrightText: © 2026 Idiap Research Institute <contact@idiap.ch>
%
% SPDX-FileContributor: Andre Freitas <andre.freitas@idiap.ch>
% Neuro-symbolic AI Group
%
% SPDX-FileContributor: Danilo Gusicuma <danilo.gusicuma@idiap.ch>
% Neuro-symbolic AI Group
%
% SPDX-License-Identifier: MIT
% Generated from step-explanations.json
\paragraph{Step-by-step reconstruction.}
A completed cycle-graph example defeats a size-independent spectral gap from connectedness alone. The task is to identify what this negative result enables.
The following is a source-guided illustrative reconstruction. Internal procedures are prescriptions; the native declarations do not establish their execution.
\begin{enumerate}[leftmargin=*,label=\arabic*.]
\item \textbf{Interpret a possible contribution} (\texttt{a0}).
\begin{description}[style=standard,labelindent=0pt,leftmargin=5.5em,labelwidth=5em,labelsep=0.5em,align=left,font=\normalfont\bfseries,itemsep=0.15em,parsep=0pt,topsep=0.3em]
\item[Purpose] Identify the bounded result’s scientific use.
\item[Input] Result-atlas slice; assumptions; failure evidence
\item[Method] Propose the smallest standalone bridge, obstruction or reframing and identify who could use it.
\item[Check] A contribution hypothesis must not inherit the breadth of the original unsolved target.
\item[Output] Contribution hypothesis: obstruction to the insufficient premise
\item[Next] Refutation tests this interpretation, not merely the underlying calculation.
\item[Example] Propose that the cycle example shows why connectedness is too weak for the intended uniform diffusion bound and motivates a stronger structural condition.
\end{description}
\item \textbf{Challenge support and significance} (\texttt{r0}).
\begin{description}[style=standard,labelindent=0pt,leftmargin=5.5em,labelwidth=5em,labelsep=0.5em,align=left,font=\normalfont\bfseries,itemsep=0.15em,parsep=0pt,topsep=0.3em]
\item[Purpose] Test contribution scope and novelty support.
\item[Input] Proposed obstruction interpretation; stated scope
\item[Method] Test hidden premises, proxy drift and prior-art boundaries; compare what the result actually rules out.
\item[Check] Neither one slow family nor a true calculation establishes a general impossibility of fast diffusion.
\item[Output] Support assessment: bounded obstruction; unestablished breadth and novelty
\item[Next] Final readiness joins this finding with the same contribution hypothesis.
\item[Example] Reject the claim that fast diffusion is impossible in general; retain only failure of connectedness to imply a uniform gap over the specified family.
\end{description}
\item \textbf{Assess value and enablement} (\texttt{e0}).
\begin{description}[style=standard,labelindent=0pt,leftmargin=5.5em,labelwidth=5em,labelsep=0.5em,align=left,font=\normalfont\bfseries,itemsep=0.15em,parsep=0pt,topsep=0.3em]
\item[Purpose] Assess target value and communication readiness.
\item[Input] Contribution hypothesis + its support assessment
\item[Method] Separate conceptual value, target alignment and actual downstream enablement; compare the prior missing condition with the next task's requirements.
\item[Check] Exploration payoff cannot by itself establish paper readiness or novelty.
\item[Output] Readiness record: useful obstruction; open merit burden
\item[Next] Original policy can route to elaboration or stop; returned audit components are not extra native occurrences.
\item[Example] The result enables replacing an inadequate premise by a search for sufficient structural assumptions; any claim of publishable novelty still requires separate evidence.
\end{description}
\end{enumerate}
\paragraph{Changed-input contrast.}
If the result merely restates a known example without a new warranted use, its mathematical validity does not establish novelty or publication merit. Enablement needs a specific downstream change.
